\documentclass[11pt,a4paper]{article}
\usepackage[T1]{fontenc}
\usepackage[utf8]{inputenc}
\usepackage[english]{babel}
\usepackage{amsmath}
\usepackage{amssymb}
\usepackage{url}
\usepackage{xcolor}
\usepackage[margin=2.5cm]{geometry}
\usepackage{hyperref}

\newcommand{\rcom}[1]{\par\noindent\textbf{\textcolor{blue!60!black}{Reviewer comment:}} \emph{#1}\par}
\newcommand{\rans}[1]{\par\noindent\textbf{\textcolor{green!40!black}{Response:}} #1\par\vspace{0.4em}}
\newcommand{\where}[1]{\par\noindent\textbf{\textcolor{orange!70!black}{Where in manuscript:}} #1\par}

\title{Response to Reviewers\\
\large Manuscript: ``Learning-Enhanced Security-Constrained Unit Commitment Problem''}
\author{Egor Grishin}
\date{}

\begin{document}
\maketitle

We thank the reviewers for their careful reading and their detailed comments. In this revision we have (i) added bootstrap 95\% confidence intervals for the main speedup summaries, (ii) added a compact ablation of the GNN screener against the 1-NN baseline and classifier-level GNN precision/recall/F1 diagnostics, (iii) added fixed-$K$ lazy-cut ablations for BANDIT, (iv) added a representative \texttt{grbtune} sanity check, (v) added a quantitative comparison to prior ML-accelerated SCUC results in the Introduction, (vi) added a dedicated \emph{Limitations and threats to validity} subsection, (vii) made the post-contingency and base-case flow constraints symmetric in notation, and (viii) clarified why GNN/GRU/BANDIT are reported as negative or marginal results. To keep the manuscript within the page limit, the largest full per-mode tables are now represented in the main paper by compact summary tables and heatmaps, with detailed CSV tables released as paper artifacts.

Below we answer each comment explicitly, in the order of the review reports. Section / equation references are to the revised manuscript.

\section*{Reviewer 1}

\rcom{R1-C1. The overall methodology of constraint reduction is not very clear. Show initial and reduced numbers of constraints.}
\rans{The revised PRUNE-$\tau$ subsection now gives the complete screening rule, including the 1-NN load feature vector, predicted injections, the emergency-margin score $s^m_\ell$, and the pruning rule $s^m_\ell\ge\tau$. The experiments section now visualizes the relationship between retained post-contingency constraints and speedup in Fig.~\ref{fig:constraint_reduction_speedup}. The supporting CSV reports the initial and retained constraint counts; for WARM+PRUNE-0.10 the reduction exceeds $99.7\%$ on the three largest systems.}
\where{Sec.~3, ``Constraint Screening (PRUNE-$\tau$)''; Sec.~\ref{sec:experiments}, Fig.~\ref{fig:constraint_reduction_speedup}; supporting CSV in \texttt{paper/data/constraint\_reduction.csv}.}

\rcom{R1-C2. The ``10x speedups'' claim in the abstract is questionable.}
\rans{The abstract no longer claims a generic $10\times$ speedup. It now reads ``up to $25.8\times$ on \texttt{case89pegase} and $9.6\times$ on average across the six benchmark cases'', matching the compact main-results table and heatmap.}
\where{Abstract, first paragraph.}

\rcom{R1-C3. No quantitative comparison with prior ML-SCUC work [10]--[13].}
\rans{We added a quantitative paragraph in the Introduction. It reports the $4.3\times$ headline number from Xavier \textit{et al.}~\cite{Xavier2019} on a 6515-bus system and the $2$--$10\times$ range of Jim\'enez-Cordero \textit{et al.}~\cite{JimenezCordero2022}, and explains that our best numbers are in the same order of magnitude but obtained from a strictly simpler pipeline that preserves exact N-1 feasibility by construction. A full head-to-head reproduction is out of scope for a benchmark-style paper, but the context is now explicit.}
\where{Introduction, paragraph starting ``Machine learning has been used to predict commitment schedules\ldots''.}

\rcom{R1-C4. The case57 ``77/23/0'' column suggests $23\%$ of instances are infeasible regardless of method.}
\rans{We clarified this in the experiments section: the same 14 dates are reported infeasible by \emph{all} 23 tested modes, which points to structural infeasibility of those MATPOWER instances (insufficient reserve after the worst-case contingency) rather than to a screening or warm-start failure. These 14 dates are excluded from all speedup statistics, so they cannot bias the acceleration factors. The exact per-date audit is stored in \texttt{paper/data/case57\_audit.csv}.}
\where{Sec.~\ref{sec:experiments}, paragraph on the robustness audit and case57 infeasibility.}

\rcom{R1-C5. GNN: missing details, no P/R/F1, and the classifier is no better than 1-NN. Please explain.}
\rans{We added the missing GNN details: line-graph construction, line-node features, criticality threshold $y_{\mathrm{thr}}=0.70$, inference threshold $\theta=0.60$, and the GraphSAGE training setup. We also added classifier-level diagnostics: on \texttt{case300}, where positive held-out labels are present at a non-negligible rate, the retrained GNN reaches precision 0.83, recall 0.51 and F1 0.63; on smaller cases, positives are absent or extremely sparse, so these metrics are not well-defined. Downstream, the GNN/1-NN speedup ratios range from 0.987 to 1.035 across PRUNE thresholds, so we explicitly treat the GNN as an exploratory negative result rather than as a recommended default.}
\where{Sec.~3, ``Constraint Screening (GNN)'' paragraph; Sec.~\ref{sec:experiments}, Tab.~\ref{tab:ablation_sanity}; supporting CSVs \texttt{paper/tables/gnn\_pr\_f1.csv} and \texttt{paper/tables/gnn\_vs\_1nn.csv}.}

\rcom{R1-C6. BANDIT: missing formula, $\mathcal{K}$, $\varepsilon$; marginal gain over LAZY.}
\rans{We now state explicitly that $\mathcal{K}=\{64,128,256\}$, the reward is $R_t=-N_t^{\mathrm{added}}$ (negative number of lazy cuts added), and $\varepsilon=0.1$. We also added fixed-$K$ ablations: static budgets can be fast on \texttt{case118}, but on \texttt{case300} all fixed-$K$ settings are slower than plain LAZY/BANDIT and pass the strict violation audit in only 66\% of runs. The text therefore presents BANDIT and fixed-$K$ cut budgeting as negative or marginal results rather than recommended techniques.}
\where{Sec.~3, ``Lazy Enforcement (LAZY, BANDIT)''; Sec.~\ref{sec:experiments}, Tab.~\ref{tab:ablation_sanity}.}

\rcom{R1-C7. GRU degrades speedup; explain why.}
\rans{We added a two-sentence diagnosis. First, MSE-trained reserve fractions $\hat r_{g,t}$ do not preserve ramp or min-up/down relationships, so the resulting dispatch violates inter-temporal feasibility more often than a direct 1-NN copy of a fully solved historical dispatch and is overwritten by the repair heuristic. Second, on small cases the solver reaches root feasibility in milliseconds, so PyTorch inference overhead is no longer negligible relative to the total solve time. The GRU is explicitly reported as a secondary ablation, not as a recommended default.}
\where{Sec.~3, ``Warm Starts (WARM, GRU)'' paragraph (revised).}

\rcom{R1-C8. Cases up to 300 buses are too small; consider larger datasets.}
\rans{We now state this limitation explicitly rather than over-claiming beyond the available benchmark. The revised paper keeps the experimental scope to the six MATPOWER cases up to \texttt{case300}; larger systems are listed as future work in the limitations section.}
\where{Sec.~\ref{sec:limitations}, \emph{Scale} bullet.}

\rcom{R1-C9. No statistical significance measures; no CI.}
\rans{We computed bootstrap 95\% confidence intervals (10\,000 resamples) for every key mode and case, and paired Wilcoxon signed-rank tests against RAW. To keep the manuscript within the page limit, the full CI table is released as a paper artifact, while the main paper reports the headline CIs inline: PRUNE-0.10 on \texttt{case89pegase} is $25.83\times$ with 95\% CI $[24.40,27.26]$, and LAZY+BANDIT on \texttt{case300} is $17.38\times$ with 95\% CI $[15.39,19.39]$. The Wilcoxon results remain significant on every benchmark case ($p<10^{-7}$).}
\where{Sec.~\ref{sec:experiments}; supporting CSVs \texttt{paper/data/speedup\_summary.csv} and \texttt{paper/data/wilcoxon.csv}.}

\rcom{R1-C10. 70/15/15 without cross-validation.}
\rans{Full $k$-fold cross-validation of the end-to-end solver pipeline is prohibitively expensive at the scale of 7\,702 solver runs: every fold would require re-solving roughly $1\,100$ MILPs (several tens of hours of compute). We therefore added this as an explicit limitation in the \emph{Limitations and threats to validity} subsection (Sec.~\ref{sec:limitations}), alongside the remark that the reported Wilcoxon tests concern per-instance differences (not split-level stability). We also state that the random seed of the 70/15/15 split is fixed and that all split indices will be released together with the code.}
\where{Sec.~\ref{sec:limitations}.}

\rcom{R1-C11. No code/data repository.}
\rans{A public release of all ML training pipelines, solver wrappers, and post-processing scripts (together with the benchmark instances and raw per-instance runtime logs) is now announced in the \emph{Reproducibility} paragraph at the end of Sec.~\ref{sec:experiments}: \url{https://github.com/nemakgregor/SCUCa}.}
\where{Sec.~\ref{sec:experiments}, \emph{Reproducibility} paragraph.}

\rcom{R1-C12. The conclusion claim is too strong.}
\rans{The opening sentence of the Conclusion now reads ``This study demonstrates that, among the methods considered in this work, the most effective way to accelerate a raw MILP formulation of SCUC is to reduce or delay post-contingency constraints\ldots''. The qualifier was present in the previous revision; we kept it and applied the same restraint to the abstract.}
\where{Conclusion, first sentence; Abstract.}

\subsection*{Editorial / typography (R1 typ-1\ldots typ-7)}
\begin{itemize}
\item typ-1 (Tab.~1 position / size): the oversized all-mode table was removed from the main manuscript and replaced by compact summary tables and released CSV artifacts.
\item typ-2 (duplicate ``This is achieved by setting\ldots''): removed in the HINTS paragraph.
\item typ-3 (Tab.~4 vs Tab.~3 labeling): the labels are now consistent.
\item typ-4 (``These issues'' $\to$ ``These observations''): fixed.
\item typ-5 (double period ``guarantees..''): fixed.
\item typ-6 (missing quotation marks in refs): fixed in \texttt{thebibliography}.
\item typ-7 (duplicate DOI in [13]): fixed.
\end{itemize}

\section*{Reviewer 2}

\rcom{R2-C1. In (12), (15), (21) the right-hand side should be $-F^{\mathrm{N}}$ / $-F^{\mathrm{E}}$.}
\rans{We rewrote the base-case and post-contingency flow bounds in symmetric two-sided form (Eqs.~\eqref{eq:base-lims}, \eqref{eq:cont-line}, \eqref{eq:cont-gen}). The lower bound $-F^{\mathrm{N}}-\underline{o}$ / $-F^{\mathrm{E}}-\underline{c}$ is now explicit; the pair of linear inequalities actually passed to the solver is stated as an equivalent form immediately after. This should match the form the reviewer expected.}
\where{Sec.~2, Eqs.~\eqref{eq:base-lims}, \eqref{eq:cont-line}, \eqref{eq:cont-gen}.}

\rcom{R2-C2. Two loads that differ only by an additive constant have $d(i)=0$ under pure z-scoring.}
\rans{The feature vector used by PRUNE-$\tau$ and by the WARM warm-start was augmented with level-sensitive summary statistics: $\phi=[z_0,\dots,z_{T-1},\mu_L,\max_t L_t,\sum_t L_t,\sigma_L]$. Two load trajectories that differ by a uniform shift now have a strictly positive distance through the $\mu_L$, $\max_t L_t$ and $\sum_t L_t$ components. The paper explicitly describes this as avoiding the degeneracy the reviewer identified.}
\where{Sec.~3, ``Constraint Screening (PRUNE-$\tau$)'' paragraph, Eq.~(20).}

\rcom{R2-C3. Could the gain be reduced to common MIP knowledge? The solver output is not discussed. Why not make (12) lazy too?}
\rans{We added an explicit paragraph in ``Lazy Enforcement (LAZY, BANDIT)'' stating why we apply lazy enforcement only to the post-contingency constraints~\eqref{eq:cont-line}--\eqref{eq:cont-gen}: they are by far the most numerous, and their threshold-style slack penalties in the objective weaken the dual bound, so pushing them out of the root LP is exactly the kind of ``threshold-reduction'' mechanism the reviewer identified. We also state that a preliminary experiment that additionally made the base-case bounds~\eqref{eq:base-lims} lazy did not yield consistent improvement. In the experiments section we added a compact sanity table with grbtune and fixed-$K$ cut-budget results, which helps separate solver tuning effects from constraint-management effects.}
\where{Sec.~3, ``Lazy Enforcement (LAZY, BANDIT)'' paragraph; Sec.~\ref{sec:experiments}, Tab.~\ref{tab:ablation_sanity}.}

\rcom{R2-C4. Gurobi has a built-in auto-tuning tool; information about its results is missing.}
\rans{We added a representative \texttt{grbtune} sanity check. On \texttt{case118}, the tuned baseline was slightly slower than default (0.93x); on \texttt{case300}, it was only 1.02x faster. These results do not replace a full per-instance tuning benchmark, but they rule out the simplest explanation that the observed order-of-magnitude gains are obtainable from Gurobi parameter tuning alone on the representative larger cases.}
\where{Sec.~\ref{sec:experiments}, Tab.~\ref{tab:ablation_sanity}; Sec.~\ref{sec:limitations}, \emph{Solver tuning} bullet.}

\vspace{1em}
\noindent We again thank the reviewers for comments that significantly sharpened the paper. We believe the revised manuscript now states its scope, evidence, and negative results accurately and defensibly.

\end{document}
